\documentclass[11pt]{article}

\usepackage[margin=1in]{geometry}
\usepackage{xcolor}
\usepackage{listings}
\usepackage{hyperref}
\usepackage{booktabs}
\usepackage{enumitem}
\usepackage{parskip}
\usepackage{titlesec}

\hypersetup{
  colorlinks = true,
  linkcolor  = blue!60!black,
  urlcolor   = blue!60!black,
  citecolor  = blue!60!black,
}

% --- Code listing style for shell commands -----------------------------------
\definecolor{shellbg}{HTML}{F5F5F5}
\definecolor{shellrule}{HTML}{CCCCCC}
\definecolor{shellcomment}{HTML}{6A737D}
\definecolor{shellkeyword}{HTML}{005CC5}

\lstdefinestyle{shell}{
  basicstyle      = \ttfamily\footnotesize,
  backgroundcolor = \color{shellbg},
  frame           = single,
  rulecolor       = \color{shellrule},
  framesep        = 6pt,
  xleftmargin     = 6pt,
  xrightmargin    = 6pt,
  breaklines      = true,
  showstringspaces= false,
  columns         = fullflexible,
  keepspaces      = true,
  commentstyle    = \color{shellcomment}\itshape,
  keywordstyle    = \color{shellkeyword}\bfseries,
  morecomment     = [l]{\#},
  morekeywords    = {sudo, apt, apt-get, yum, dnf, brew, git, make, cd, mkdir,
                     export, source, ./configure, configure, echo, tar, wget,
                     curl, ldconfig, which, ln, cp, rm, mv},
  literate        = {/}{{/}}1 {-}{{-}}1 {=}{{=}}1,
}
\lstset{style=shell}

\titleformat{\section}{\Large\bfseries}{\thesection}{0.7em}{}
\titleformat{\subsection}{\large\bfseries}{\thesubsection}{0.6em}{}

\title{Installing Magic, Netgen, IRSIM, ngspice, xschem, and open\_pdks\\
       \large An installation guide for the Open Circuit Design tool suite}
\author{James E. Stine}
\date{\today}

\begin{document}
\maketitle

\section{Overview}

This guide walks through building and installing the open-source analog/digital
IC design flow maintained primarily by Tim Edwards into a self-contained user
prefix. The default location used throughout this document is

\begin{lstlisting}
$HOME/opencircuitdesign
\end{lstlisting}

\noindent which matches the default \texttt{OPEN\_CIRCUITDESIGN\_ROOT} value
in the included \texttt{setup.sh}. If you would rather install elsewhere
(e.g.\ \texttt{/opt/opencircuitdesign} or \texttt{/home/user/ocd}), simply
edit the line near the top of \texttt{setup.sh}:

\begin{lstlisting}
export OPEN_CIRCUITDESIGN_ROOT="$HOME/opencircuitdesign"
\end{lstlisting}

\noindent and use that same path everywhere this guide writes
\texttt{\$OPEN\_CIRCUITDESIGN\_ROOT}.

The tools covered are:

\begin{itemize}[leftmargin=2em]
  \item \textbf{Magic} --- VLSI layout editor (\url{http://opencircuitdesign.com/magic}).
  \item \textbf{Netgen} --- netlist comparison / LVS
        (\url{http://opencircuitdesign.com/netgen}).
  \item \textbf{IRSIM} --- switch-level digital simulator
        (\url{http://opencircuitdesign.com/irsim}).
  \item \textbf{ngspice} --- analog/mixed-signal SPICE simulator
        (\url{https://ngspice.sourceforge.io}). \emph{Note:} ngspice is a
        separate project from Tim Edwards' tools, but is required by
        \texttt{open\_pdks} and is conventionally bundled with the suite.
  \item \textbf{xschem} --- schematic capture
        (\url{https://github.com/StefanSchippers/xschem}). Also a separate
        upstream, included because \texttt{setup.sh} configures it.
  \item \textbf{open\_pdks} --- installer/wrapper for open-source PDKs such
        as \texttt{sky130A} and \texttt{gf180mcu}
        (\url{https://github.com/RTimothyEdwards/open_pdks}).
\end{itemize}

Everything will be installed under a single prefix:

\begin{lstlisting}
$OPEN_CIRCUITDESIGN_ROOT
|-- bin/        # magic, netgen, irsim, ngspice, xschem executables
|-- lib/        # shared libraries
|-- share/      # PDKs, scripts, technology files, magic/netgen/xschem libs
|-- include/    # development headers
\end{lstlisting}

\section{Prerequisites}

\subsection{Build dependencies (Debian / Ubuntu)}

Install the toolchain and X/Tcl libraries before building any of the tools.

\begin{lstlisting}
sudo apt-get update
sudo apt-get install -y \
    build-essential m4 tcsh csh git \
    tcl-dev tk-dev \
    libcairo2-dev \
    mesa-common-dev libglu1-mesa-dev \
    libncurses-dev \
    libx11-dev libxpm-dev libxt-dev \
    libreadline-dev \
    bison flex \
    autoconf automake libtool \
    libfftw3-dev libxaw7-dev \
    python3 python3-pip
\end{lstlisting}

\subsection{Build dependencies (Red Hat / Fedora)}

\begin{lstlisting}
sudo dnf install -y \
    gcc gcc-c++ make m4 tcsh git \
    tcl-devel tk-devel \
    cairo-devel \
    mesa-libGL-devel mesa-libGLU-devel \
    ncurses-devel \
    libX11-devel libXpm-devel libXt-devel libXaw-devel \
    readline-devel \
    bison flex \
    autoconf automake libtool \
    fftw-devel \
    python3 python3-pip
\end{lstlisting}

\subsection{Create the install prefix and source directories}

Pick the install root once and use it throughout. The default in
\texttt{setup.sh} is \texttt{\$HOME/opencircuitdesign}; set it in your shell
\emph{first} so the build commands below pick it up:

\begin{lstlisting}
# Match the value in setup.sh (edit setup.sh first if you want a different path)
export OPEN_CIRCUITDESIGN_ROOT="$HOME/opencircuitdesign"

mkdir -p "$OPEN_CIRCUITDESIGN_ROOT"
mkdir -p "$OPEN_CIRCUITDESIGN_ROOT/src"   # all sources will be cloned here
cd "$OPEN_CIRCUITDESIGN_ROOT/src"
\end{lstlisting}

\subsection{Use the included \texttt{setup.sh}}

Rather than hand-rolling \texttt{PATH}, \texttt{LD\_LIBRARY\_PATH},
\texttt{PDK\_ROOT}, and the per-tool environment variables, the repository
ships a \texttt{setup.sh} script that handles all of it. After confirming
the \texttt{OPEN\_CIRCUITDESIGN\_ROOT} line at the top of \texttt{setup.sh}
matches the prefix you chose above, source the script in any shell where
you want the tools available:

\begin{lstlisting}
# One-off: source it in the current shell
source /path/to/repo/setup.sh

# Persistent: have every login shell source it automatically
echo 'source /path/to/repo/setup.sh' >> ~/.bashrc
\end{lstlisting}

\noindent The script:

\begin{itemize}[leftmargin=2em]
  \item Verifies \texttt{\$OPEN\_CIRCUITDESIGN\_ROOT} exists and bails out
        loudly if it doesn't.
  \item Prepends \texttt{\$OPEN\_CIRCUITDESIGN\_ROOT/bin} to \texttt{PATH}
        and \texttt{\$OPEN\_CIRCUITDESIGN\_ROOT/lib} to
        \texttt{LD\_LIBRARY\_PATH}.
  \item Sets \texttt{PDK\_ROOT} to
        \texttt{\$OPEN\_CIRCUITDESIGN\_ROOT/share/pdk} (the standard
        \texttt{open\_pdks} install location when configured with
        \texttt{-{}-with-pdk\_root=\$OPEN\_CIRCUITDESIGN\_ROOT/share/pdk}).
  \item Exports per-tool variables: \texttt{MAGIC\_HOME}, \texttt{NETGEN\_HOME},
        \texttt{NGSPICE\_SHARE\_DIR}, \texttt{XSCHEM\_SYSTEM\_LIBRARY\_PATH}.
  \item If \texttt{sky130A} is found under \texttt{\$PDK\_ROOT}, it also
        sets \texttt{PDK=sky130A}, \texttt{SKY130A},
        \texttt{MAGIC\_MAGICRC}, \texttt{NETGEN\_SETUP},
        \texttt{XSCHEM\_LIBRARY\_PATH}, and \texttt{XSCHEM\_RC}.
  \item Prints a short summary so you can confirm everything resolved
        correctly.
\end{itemize}

You should source \texttt{setup.sh} \emph{before} running the verification
commands at the end of each section below. The build steps themselves do
not need it --- they only need \texttt{OPEN\_CIRCUITDESIGN\_ROOT} exported.

\section{Magic}

Magic is the VLSI layout editor used as the geometry engine for the entire flow.

\begin{lstlisting}
cd "$OPEN_CIRCUITDESIGN_ROOT/src"
git clone https://github.com/RTimothyEdwards/magic.git
cd magic
./configure --prefix="$OPEN_CIRCUITDESIGN_ROOT"
make -j$(nproc)
make install
\end{lstlisting}

Verify (after \texttt{source setup.sh}):

\begin{lstlisting}
which magic
magic -dnull -noconsole -T scmos <<< "version; quit -noprompt"
\end{lstlisting}

\section{Netgen}

Netgen is the LVS (Layout-Versus-Schematic) and netlist comparison tool.

\begin{lstlisting}
cd "$OPEN_CIRCUITDESIGN_ROOT/src"
git clone https://github.com/RTimothyEdwards/netgen.git
cd netgen
./configure --prefix="$OPEN_CIRCUITDESIGN_ROOT"
make -j$(nproc)
make install
\end{lstlisting}

Verify:

\begin{lstlisting}
which netgen
netgen -batch <<< "puts [version]; quit"
\end{lstlisting}

\section{IRSIM}

IRSIM is a switch-level simulator that integrates with Magic.

\begin{lstlisting}
cd "$OPEN_CIRCUITDESIGN_ROOT/src"
git clone https://github.com/RTimothyEdwards/irsim.git
cd irsim
./configure --prefix="$OPEN_CIRCUITDESIGN_ROOT" \
            --with-tcl=/usr/lib/tcl8.6 \
            --with-tk=/usr/lib/tk8.6
make -j$(nproc)
make install
\end{lstlisting}

If your distribution uses a different Tcl/Tk version, adjust the
\texttt{-{}-with-tcl}/\texttt{-{}-with-tk} flags accordingly (or omit them to
let \texttt{configure} auto-detect).

Verify:

\begin{lstlisting}
which irsim
irsim -V
\end{lstlisting}

\section{ngspice}

ngspice is required by \texttt{open\_pdks} for simulation models and
verification. It is built from a release tarball rather than git.

\begin{lstlisting}
cd "$OPEN_CIRCUITDESIGN_ROOT/src"
NGSPICE_VER=44
wget https://sourceforge.net/projects/ngspice/files/ng-spice-rework/${NGSPICE_VER}/ngspice-${NGSPICE_VER}.tar.gz
tar xzf ngspice-${NGSPICE_VER}.tar.gz
cd ngspice-${NGSPICE_VER}
mkdir -p release && cd release
../configure --prefix="$OPEN_CIRCUITDESIGN_ROOT" \
             --enable-xspice \
             --enable-cider \
             --enable-openmp \
             --with-readline=yes \
             --disable-debug
make -j$(nproc)
make install
\end{lstlisting}

\noindent\textbf{Notes:}
\begin{itemize}[leftmargin=2em]
  \item Adjust \texttt{NGSPICE\_VER} to the latest stable release; check
        \url{https://ngspice.sourceforge.io/download.html}.
  \item \texttt{-{}-enable-xspice} is required for code-model support used by
        many open-source PDKs.
  \item Building inside a \texttt{release/} subdirectory keeps the source tree
        clean as recommended by the ngspice manual.
\end{itemize}

Verify:

\begin{lstlisting}
which ngspice
ngspice -v
\end{lstlisting}

\section{xschem}

\texttt{setup.sh} configures xschem variables
(\texttt{XSCHEM\_SYSTEM\_LIBRARY\_PATH}, \texttt{XSCHEM\_LIBRARY\_PATH},
\texttt{XSCHEM\_RC}), so install it into the same prefix.

\begin{lstlisting}
cd "$OPEN_CIRCUITDESIGN_ROOT/src"
git clone https://github.com/StefanSchippers/xschem.git
cd xschem
./configure --prefix="$OPEN_CIRCUITDESIGN_ROOT"
make -j$(nproc)
make install
\end{lstlisting}

Verify:

\begin{lstlisting}
which xschem
xschem --version
\end{lstlisting}

\section{open\_pdks}

\texttt{open\_pdks} is a meta-installer: it fetches and configures vendor PDKs
(sky130, gf180mcu, etc.) and stitches them into Magic, Netgen, ngspice, and
xschem. It must be configured \emph{after} the tools above are on
\texttt{PATH} so it can locate them.

\subsection{Configure and build}

Note the underscore in \texttt{-{}-with-pdk\_root}: this is what
\texttt{setup.sh} expects when it sets
\texttt{PDK\_ROOT="\$OPEN\_CIRCUITDESIGN\_ROOT/share/pdk"}.

\begin{lstlisting}
cd "$OPEN_CIRCUITDESIGN_ROOT/src"
git clone https://github.com/RTimothyEdwards/open_pdks.git
cd open_pdks

./configure --prefix="$OPEN_CIRCUITDESIGN_ROOT" \
            --with-pdk_root="$OPEN_CIRCUITDESIGN_ROOT/share/pdk" \
            --enable-sky130-pdk \
            --enable-gf180mcu-pdk \
            --with-magic="$OPEN_CIRCUITDESIGN_ROOT/bin/magic" \
            --with-netgen="$OPEN_CIRCUITDESIGN_ROOT/bin/netgen" \
            --with-ngspice="$OPEN_CIRCUITDESIGN_ROOT/bin/ngspice"
\end{lstlisting}

The \texttt{--enable-*-pdk} flags select which PDKs to fetch. Drop any you
don't need; each PDK is several gigabytes after install.

\subsection{Install (this step is long --- expect 30--90 minutes)}

\begin{lstlisting}
make
make install
\end{lstlisting}

The PDKs land under \texttt{\$PDK\_ROOT}, e.g.
\texttt{\$OPEN\_CIRCUITDESIGN\_ROOT/share/pdk/sky130A}. Once
\texttt{sky130A} exists, sourcing \texttt{setup.sh} again will pick it up
and export \texttt{PDK}, \texttt{SKY130A}, \texttt{MAGIC\_MAGICRC},
\texttt{NETGEN\_SETUP}, \texttt{XSCHEM\_LIBRARY\_PATH}, and
\texttt{XSCHEM\_RC} automatically.

\subsection{Verify}

\begin{lstlisting}
source /path/to/repo/setup.sh

ls "$PDK_ROOT"
# Expect: sky130A  sky130B  gf180mcuA  gf180mcuB  gf180mcuC  gf180mcuD  ...

# setup.sh should now report sky130A as the selected PDK:
echo "$PDK"           # -> sky130A
echo "$MAGIC_MAGICRC" # -> .../sky130A/libs.tech/magic/sky130A.magicrc
echo "$NETGEN_SETUP"  # -> .../sky130A/libs.tech/netgen/sky130A_setup.tcl

# Open a sky130 layout in magic to confirm tech files were installed:
magic -rcfile "$MAGIC_MAGICRC"
\end{lstlisting}

\section{Putting it all together}

\subsection{Quick sanity check}

After every tool is installed and you have sourced \texttt{setup.sh},
the script's own summary output should list each tool's location. You can
also run:

\begin{lstlisting}
for tool in magic netgen irsim ngspice xschem; do
    printf "%-8s -> %s\n" "$tool" "$(which $tool 2>/dev/null || echo '<not found>')"
done
\end{lstlisting}

All paths should resolve under \texttt{\$OPEN\_CIRCUITDESIGN\_ROOT/bin}.

\subsection{Recommended directory layout after install}

\begin{lstlisting}
$OPEN_CIRCUITDESIGN_ROOT
|-- bin/
|   |-- magic
|   |-- netgen
|   |-- irsim
|   |-- ngspice
|   `-- xschem
|-- include/
|-- lib/
|-- share/
|   |-- magic/      # MAGIC_HOME
|   |-- netgen/     # NETGEN_HOME
|   |-- ngspice/    # NGSPICE_SHARE_DIR
|   |-- xschem/     # XSCHEM_SYSTEM_LIBRARY_PATH
|   |-- pdk/        # PDK_ROOT
|   |   |-- sky130A/
|   |   `-- gf180mcuA/
|   `-- man/
`-- src/            # cloned source trees
    |-- magic/
    |-- netgen/
    |-- irsim/
    |-- ngspice-44/
    |-- xschem/
    `-- open_pdks/
\end{lstlisting}

\section{Troubleshooting}

\begin{description}[leftmargin=0em, style=nextline]
  \item[\texttt{ERROR: Directory \$OPEN\_CIRCUITDESIGN\_ROOT does not exist}]
    \texttt{setup.sh} aborts when its target prefix is missing. Either
    create the directory (\texttt{mkdir -p "\$OPEN\_CIRCUITDESIGN\_ROOT"})
    or edit the \texttt{OPEN\_CIRCUITDESIGN\_ROOT=...} line at the top of
    \texttt{setup.sh} to point somewhere that does exist.

  \item[\texttt{configure: error: Tcl headers not found}]
    Make sure \texttt{tcl-dev} / \texttt{tcl-devel} is installed and that
    \texttt{tclConfig.sh} lives in a standard location (typically
    \texttt{/usr/lib/tclX.Y/}). Pass \texttt{-{}-with-tcl=<dir>} explicitly
    if needed.

  \item[Magic launches but graphics are missing]
    Build Magic with Cairo and OpenGL support enabled (the default when
    the corresponding \texttt{-dev} packages are present). Re-run
    \texttt{./configure} after installing them and rebuild.

  \item[\texttt{open\_pdks} fails to find Magic / Netgen]
    Ensure \texttt{\$OPEN\_CIRCUITDESIGN\_ROOT/bin} is on \texttt{PATH}
    before running \texttt{./configure}, or pass the absolute paths via
    \texttt{-{}-with-magic=}, \texttt{-{}-with-netgen=},
    \texttt{-{}-with-ngspice=} as shown above.

  \item[\texttt{setup.sh} warns ``\texttt{No PDK directory found}'']
    \texttt{open\_pdks} hasn't been built/installed yet, or it was
    configured with a different \texttt{-{}-with-pdk\_root}. Re-run
    \texttt{./configure} for \texttt{open\_pdks} with
    \texttt{-{}-with-pdk\_root="\$OPEN\_CIRCUITDESIGN\_ROOT/share/pdk"}.

  \item[Shared libraries not found at runtime]
    Source \texttt{setup.sh} (it sets \texttt{LD\_LIBRARY\_PATH}). For a
    system-wide alternative, add \texttt{\$OPEN\_CIRCUITDESIGN\_ROOT/lib}
    to \texttt{/etc/ld.so.conf.d/opencircuitdesign.conf} and run
    \texttt{sudo ldconfig}.

  \item[Updating later]
    Each tool can be rebuilt independently:
    \texttt{cd \$OPEN\_CIRCUITDESIGN\_ROOT/src/<tool> \&\& git pull \&\& make \&\& make install}.
    Re-run \texttt{open\_pdks} only when changing PDK selection or after a
    significant Magic/Netgen update.
\end{description}

\section{References}

\begin{itemize}[leftmargin=2em]
  \item Open Circuit Design home page: \url{http://opencircuitdesign.com}
  \item Magic on GitHub: \url{https://github.com/RTimothyEdwards/magic}
  \item Netgen on GitHub: \url{https://github.com/RTimothyEdwards/netgen}
  \item IRSIM on GitHub: \url{https://github.com/RTimothyEdwards/irsim}
  \item ngspice: \url{https://ngspice.sourceforge.io}
  \item xschem on GitHub: \url{https://github.com/StefanSchippers/xschem}
  \item open\_pdks on GitHub: \url{https://github.com/RTimothyEdwards/open_pdks}
\end{itemize}

\end{document}
